\subsection{谱投影}

设 E 为 Banach 空间。我们用 A 表示由 E 的自同态组成的幺 Banach 代数 \(\mathcal{L}(\mathrm{E})\) 。设 \(u\in \mathrm{A}\) 。

设 H 是 \(\mathrm{Sp}(u)\) 的一个在 \(\mathrm{Sp}(u)\) 中既开又闭的子集；用 K 表示它在 \(\mathrm{Sp}(u)\) 中的补集。

与 \(u\) 和 H 相伴的幂等元（I，第 81 页，第 12 节）是 E 的一个连续投影，称为与 \(u\) 和 H 相伴的谱投影，记作 \(e_{\mathrm{H}}(u)\) ，或简记为 \(e_{\mathrm{H}}\) 。它的像称为 E 的与 \(u\) 和 H 相伴的谱子空间，记作 \(\mathrm{E_H}(u)\) ，或简记为 \(\mathrm{E_H}\) 。它的核是 \(\mathrm{E_K}(u)\) ，也记作 \(\widetilde{\mathrm{E}}_{\mathrm{H}}(u)\) ，或简记为 \(\widetilde{\mathrm{E}}_{\mathrm{H}}\) 。空间 E 是闭子空间 \(\mathrm{E_H}\) 与 \(\mathrm{E_K}\) 的拓扑直和。条件 \(\mathrm{E_H} = 0\) 等价于 \(e_{\mathrm{H}} = 0\) ，即等价于 H 的指示函数 \(f_{\mathrm{H}}\) 在 \(\mathrm{Sp}(u)\) 上为零函数，从而等价于 \(\mathrm{H} = \varnothing\) 。

每个与 u 交换的 E 的自同态 v 也与 \(e_{\mathrm{H}}(u)\) 交换（I，第 74 页，定理 5），因此它使 \(E_{H}\) 与 \(E_{K}\) 保持稳定。特别地，自同态 u 使子空间 \(E_{H}\) 与 \(E_{K}\) 保持稳定。空间 \(E_{K}\) 是 \(E_{H}\) 在 E 中的唯一的在 u 下稳定的拓扑补。

幺代数 \(A_{H} = e_{H}Ae_{H}\) （同前）由 E 的使 \(E_{H}\) 保持稳定且在 \(E_{K}\) 上为零的那些自同态组成，它是 A 的子代数。对每个 \(v \in A_{H}\) ，用 \(v|E_{H}\) 表示从 v 诱导出的 \(E_{H}\) 的自同态。映射 \(v \mapsto v|E_{H}\) 是从 \(A_{H}\) 到 \(\mathcal{L}(E_{\mathrm{H}})\) 上的同构。特别地，由 I，第 82 页，公式 (18)，我们有 \(\mathrm{Sp}(u|E_{\mathrm{H}}) = \mathrm{Sp}_{\mathrm{A}_{\mathrm{H}}} (u|E_{\mathrm{H}}) = H\) 及 \(\mathrm{Sp}_{\mathrm{A}_{\mathrm{K}}} (u|E_{\mathrm{K}}) = K\) 。

命题 2.　设 E 为 Banach 空间，u 为 E 的自同态。设 \(E_{1}\) 与 \(E_{2}\) 是 E 的在 u 下稳定的闭子空间，满足 \(E = E_{1} \oplus E_{2}\) 。假设自同态 \(u_{1}\) 与 \(u_{2}\) （即从 u 诱导出的 \(E_{1}\) 与 \(E_{2}\) 的自同态）具有不相交的谱 \(H_{1} = \mathrm{Sp}(u_{1})\)

与 \(\mathrm{H}_2 = \mathrm{Sp}(u_2)\) 。则 \(\mathrm{Sp}(u) = \mathrm{H}_1 \cup \mathrm{H}_2\) ，且 \(e_{\mathrm{H}_1}(u)\) 是以 \(\mathrm{E}_1\) 为像、以 \(\mathrm{E}_2\) 为核的投影。特别地，我们有 \(\mathrm{E}_{\mathrm{H}_1} = \mathrm{E}_1\) 与 \(\mathrm{E}_{\mathrm{H}_2} = \mathrm{E}_2\) 。

我们有 \(\mathrm{Sp}(u)=\mathrm{H}_{1}\cup\mathrm{H}_{2}\) （I，第 128 页，引理 1）。由于集合 \(H_{1}\) 与 \(H_{2}\) 是紧的，它们在 \(\mathrm{Sp}(u)\) 中既开又闭。对每个在 \(\mathrm{Sp}(u)\) 的邻域中全纯的函数 f，自同态 \(f(u)\) 使 \(E_{1}\) 与 \(E_{2}\) 保持稳定，且在 \(E_{1}\) 上与 \(f(u_{1})\) 一致，在 \(E_{2}\) 上与 \(f(u_{2})\) 一致（同前）。特别地，取 f 为全纯函数的一个芽 \(f_{H_{1}}\) ，它在 \(H_{1}\) 的邻域中等于 1，在 \(H_{2}\) 的邻域中等于 0（参见 I，第 81 页，第 12 节）；则 \(f_{\mathrm{H}_{1}}(u_{1})\) 是 \(E_{1}\) 的恒等映射，因为 \(f_{H_{1}}=1\) 在 \(\mathrm{H}_{1}=\mathrm{Sp}(u_{1})\) 的某个邻域中成立；而 \(f_{\mathrm{H}_{1}}(u_{2})\) 为零，因为 \(f_{H_{1}}=0\) 在 \(H_{2}\) 的某个邻域中成立。因此 \(e_{\mathrm{H}_{1}}(u)=f_{\mathrm{H}_{1}}(u)\) 是以 \(E_{1}\) 为像、以 \(E_{2}\) 为核的投影，于是我们有 \(E_{1}=E_{H_{1}}\) 与 \(E_{2}=E_{H_{2}}\) 。

设 E 为 Banach 空间，u 为 E 的自同态。设 \((\mathrm{H}_{i})_{i\in\mathrm{I}}\) 为 \(\mathrm{Sp}(u)\) 的两两不相交的既开又闭子集的有限族，H 为它们的并。I，第 81 页，关系式 (15) 与 (16) 蕴涵下列断言：

\begin{enumerate}
\def\labelenumi{\alph{enumi})}
\item
  投影族 \((e_{\mathrm{H}_{i}}(u))_{i\in\mathrm{I}}\) 是正交的（即 \(e_{\mathrm{H}_{i}}(u)e_{\mathrm{H}_{j}}(u)=0\) 对所有 \((i,j)\) ∈ I²（\(i\neq j\)）成立，参见 A，II，第 18 页，定义 7），且其和为 \(e_{\mathrm{H}}(u)\) ；
\item
  向量空间 \(E_{H}\) 是族 \((\mathrm{E}_{\mathrm{H}_{i}})_{i\in\mathrm{I}}\) 的拓扑直和；
\item
  对每个 \(j \in \mathrm{I}\) ，有 \(\mathrm{E}_{\mathrm{Sp}(u) - \mathrm{H}_j} = \mathrm{E}_{\mathrm{Sp}(u) - \mathrm{H}} \oplus \bigoplus_{i \neq j} \mathrm{E}_{\mathrm{H}_i}\) ；
\item
  向量空间 \(\mathrm{E}_{\mathrm{Sp}(u)-\mathrm{H}}\) 是族 \((\mathrm{E}_{\mathrm{Sp}(u)-\mathrm{H}_i})_{i\in\mathrm{I}}\) 的交。
\end{enumerate}

当 \(\mathrm{H} = \mathrm{Sp}(u)\) 时，E 的拓扑直和分解 \(\bigoplus_{i \in I} E_{H_i}\) 称为 E 的与 u 及有限分割 \((\mathrm{H}_i)_{i \in I}\) （\(\mathrm{Sp}(u)\) 的一个分割）相伴的谱分解。

设 \(f\) 为 \(\mathcal{O}(\mathrm{Sp}(u))\) 中的元素。E 的自同态 \(f(u)\) 的谱是 \(f\) 作用于 \(u\) 的谱所得的像（I，第 75 页，命题 8）。设 L 是 \(\mathrm{Sp}(f(u))\) 的一个既开又闭的子集。集合 H 由这样的元素 \(\lambda \in \mathrm{Sp}(u)\) 组成：\(f(\lambda)\) 属于 L；它在 \(\mathrm{Sp}(u)\) 中既开又闭，且我们有 \(e_{\mathrm{L}}(f(u)) = e_{\mathrm{H}}(u)\) ，因为有 \(f_{\mathrm{L}} \circ f = f_{\mathrm{H}}\) 在 \(\mathcal{O}(\mathrm{Sp}(u))\) 中成立（同前）。

